\section{Related Work}\label{sec:related_work}

\paragraph{VLA robustness.} VLAs have progressed from tokenized prediction~\cite{brohan2023rt2,kim2024openvla,ghosh2024octo} to flow matching~\cite{black2024pi0,shukor2025smolvla}, yet lose over 65\,pp under deployment perturbations~\cite{fei2025liberoplus}. Training-time defenses~\cite{srpo2025,robustvla2025,digflow2025} and test-time RL adaptation~\cite{liu2026ttvla} improve robustness but require modifying VLA weights.

\paragraph{Inference-time correction.} Among frozen-VLA methods, ReconVLA~\cite{chen2026reconvla} quantifies uncertainty without learned correction; SC-VLA~\cite{liu2026scvla} trains auxiliary world-model heads for online refinement via reward reshaping; VLA-Pilot~\cite{li2025vlapilot} uses evolutionary search at 2.41\,s per step. Tang et al.~\cite{tang2026shifting} propose uncertainty quantification for VLAs via max-pooled sliding windows but do not learn corrective actions. A2C2~\cite{a2c2_2025}, the closest concurrent work, learns a correction head for temporal staleness but applies it unconditionally. Our pipeline detects \emph{when correction would hurt}, using calibrated gating to suppress correction on confident steps. Residual policy learning~\cite{silver2018residual,johannink2019residual} shares the additive structure; magnitude control is addressed by gating~\cite{jiang2024transic}, clipping~\cite{xiao2026pld}, or decay~\cite{pan2025drc}. Our CRM uses fixed $L_2$ clipping with supervised loss, requiring no RL.

\paragraph{Selective prediction.} Our $\tau$-gating is an instance of selective prediction~\cite{chow1970optimum,elyaniv2010selective,geifman2017selective}. Unlike SelectiveNet~\cite{geifman2019selectivenet}, we use a pairwise distance threshold with no additional learned components. The $K{=}2$ $L_2$ distance is a minimal ensemble disagreement measure~\cite{lakshminarayanan2017simple}, exploiting the inherent stochasticity of flow-matching generation rather than injected noise as in MC Dropout~\cite{gal2016dropout}, and requiring no architectural changes. Conformal prediction~\cite{angelopoulos2023gentle} has been applied to robotics~\cite{lindemann2023safe,zhao2024conformalized,yao2025cpnav}; our threshold draws on this calibration philosophy without formal coverage guarantees.

\paragraph{Inference-time scaling.} Test-time scaling in LLMs~\cite{snell2024scaling} uses best-of-N selection; RoVer~\cite{dai2025rover} and DA-SIP~\cite{chun2025dasip} extend this to VLAs. Our analysis reveals a different pattern: learned correction is harmful when unconditional, and the key mechanism is calibrated gating rather than scaling sample count (Section~\ref{sec:main_results}). To our knowledge, no prior work systematically examines \emph{when} post-hoc correction of frozen VLAs degrades performance or provides a calibrated mechanism to selectively suppress it. We address both the analysis---establishing that unconditional correction is universally harmful across tasks and model scales---and the method: adaptive gating that makes learned correction safe to deploy.
